\vspace{0.5em}\noindent\textbf{Event Name}\par

FCAJ x Agentic AI Build Week 2026: Show Up. Build. Pitch. WIN!

This program presented a recap of the notable activities, products, and
project presentations from Agentic AI Build Week 2026. The main program
was organized over five days, from July 8 to July 12, 2026, and included
workshops, product development activities, a hackathon, and AI solution
pitches before a panel of judges.

\vspace{0.5em}\noindent\textbf{Time}\par

\begin{itemize}
\tightlist
\item
  Agentic AI Build Week: July 8--12, 2026.
\item
  Team offline participation: July 26, 2026.
\end{itemize}

\vspace{0.5em}\noindent\textbf{Location}\par

\begin{itemize}
\tightlist
\item
  Participation format: Offline.
\item
  Location of the main program: Ho Chi Minh City, Vietnam.
\end{itemize}

\vspace{0.5em}\noindent\textbf{Role}\par

Offline participant.

During the event, our team focused on how the participating teams
identified real-world problems, developed Agentic AI solutions, selected
appropriate AWS architectures, and presented the value of their products
to the judges.

\vspace{0.5em}\noindent\textbf{Main Content}\par

The event focused on the development of AI Agent and Agentic AI systems
capable of completing multi-step tasks, using external tools, and
supporting the resolution of practical business problems.

Unlike conventional chatbots that mainly respond to questions, AI Agents
can analyze requests, create plans, select appropriate tools, process
data, and perform actions based on predefined objectives.

Several notable topics and projects were introduced during the program:

\begin{itemize}
\tightlist
\item
  Developing a multi-channel food-ordering chatbot through Zalo and
  Messenger.
\item
  Building a Multi-Agent system to collect and analyze signals related
  to changes in corporate strategies.
\item
  Automating AWS architecture design, system diagram generation, and
  deployment cost estimation.
\item
  Combining Computer Vision and Agentic AI to monitor crowd density,
  predict congestion, and support crowd coordination.
\item
  Using multiple AI Sub-agent to assist in detecting and investigating
  suspicious financial transactions.
\end{itemize}

Through the presentations, I learned that an AI product requires more
than a capable model. It must also address the correct user problem. The
teams had to balance technical architecture, development time,
operational cost, security, and business value.

Another important topic was the role of Human-in-the-loop mechanisms. In
sensitive areas such as finance, banking, and customer service, AI can
collect information, analyze data, and provide recommendations, but a
human should still review the results and make the final decision.

\vspace{0.5em}\noindent\textbf{Participation Evidence}\par

Team offline participation evidence:

\vspace{0.5em}\noindent\textbf{Lessons Learned}\par

After watching the event, I identified several important lessons.

First, an AI project should begin with a real pain point rather than
with the selection of a trending technology. A complex architecture
provides limited value if the product does not address a specific user
need.

Second, development teams should control the project scope carefully
when working under time constraints. A stable MVP that clearly
demonstrates its value is often more effective than a system with many
incomplete features.

Third, when developing Agentic AI systems, teams must consider accuracy,
security, tool permissions, model usage costs, and the ability to track
the actions performed by the Agent. AI should not be allowed to execute
critical actions without appropriate control and confirmation
mechanisms.

\vspace{0.5em}\noindent\textbf{Personal Contribution}\par

My personal contribution was to actively follow the presentations,
record the main points, and summarize the Agentic AI application
patterns introduced during the event.

I also analyzed how these concepts could be applied to my internship
project, particularly in using AI Agents to process data, support users,
and automate repetitive tasks. These notes provided additional
references for evaluating feasibility, project scope, and security
requirements when integrating AI into a real system.
